% ---------------------------------------------------------------------
%  Abstract  --  unnumbered, but listed in the table of contents
% ---------------------------------------------------------------------
\chapter*{Abstract}
\addcontentsline{toc}{chapter}{Abstract}
\markboth{Abstract}{}

\begingroup
\setlength{\leftskip}{0pt}\noindent
Collaborative writing platforms rebuild a document every time an author
pauses typing, which makes compilation latency the single most visible
performance characteristic of the system. This report presents a
reproducible benchmark of three build strategies -- a cold baseline, a
warm package cache, and incremental recompilation -- measured across
documents of \numrange{20}{160} pages. We find that incremental
recompilation keeps median build time below \qty{1.4}{\second} even for
the largest documents, an order of magnitude faster than the cold
baseline, and that the crossover point where caching stops paying off
occurs at roughly 60 pages. We close with a short discussion of what
these numbers mean for multi-author projects and how the measurement
harness can be reused.
\endgroup

\vspace{6mm}
\noindent\textbf{\sffamily Keywords:}\enspace
compilation latency \textperiodcentered\ reproducible builds
\textperiodcentered\ collaborative editing \textperiodcentered\
document engineering
